\section{Introduction}\label{sec:intro}

The Beal conjecture asserts that if $A^x+B^y=C^z$ with integers
$A,B,C\neq 0$ and exponents $x,y,z>2$, then $A,B,C$ share a prime
factor~\cite{mauldin1997beal}. This draft does not prove the
conjecture, does not address the unrestricted signature $(4,4,13)$
equation
\[
  A^4+B^4=C^{13},
\]
and does not derive any reduction of that equation. It records a
\emph{structural verification} of Lean~4.12 material attached to a
single, deliberately restricted one-parameter family inside that
signature, together with an honest account of what each Lean name
does and does not establish.

\subsection*{The gap-3 family: what it is and where it comes from}

Throughout, ``the gap-3 equation'' means
\begin{equation}\label{eq:gap3}
  A^4+B^4=(B+3)^{13},
\end{equation}
i.e.\ $C:=B+3$. This is \emph{not} a reduction proved from the
unrestricted equation anywhere in the repositories accompanying
this draft. It is the name given, inside a much larger and mostly
uninhabited Lean development (files
\texttt{BealGap3BakerUpperBound.lean}, \texttt{BealFreyB14.lean},
\texttt{DarmonMerelFrey4413.lean}, and the
\texttt{Beal\_4\_13\_13\_Zsigmondy\_13\_Step*} series), to one
specific test family that the author chose to formalize as a
vehicle for exercising Tate/Mazur/Ribet/Kolyvagin-style arguments
in Lean. Four points, requested explicitly by an external
reviewer, are recorded here rather than left implicit:
\begin{enumerate}
\item \emph{Where $C=B+3$ comes from.} It is an adopted
  normalization, not a derived one. The surrounding Lean files
  contain a size lemma for the \emph{different} triple $(4,13,13)$
  showing $C^{13}-B^{13}\ge 26B^{12}$ when $C\ge B+2$
  (\texttt{Beal\_4\_13\_13\_Size\_C\_ge\_B\_plus\_2.lean}) and an
  observation that $\gcd(B,B+3)\mid 3$
  (\texttt{gcd\_B\_Bplus3\_dvd\_three} in
  \texttt{BealFreyB14.lean}), but neither of these, nor anything
  else found in the repository, is a theorem that a hypothetical
  $(4,4,13)$ counterexample must have $C-B=3$. The gap-3 family is
  a case study, not a WLOG reduction.
\item \emph{What it covers.} Only the one-parameter slice
  \eqref{eq:gap3}, for arbitrary $A,B\in\mathbb N$. It says nothing
  about $C-B\in\{1,2,4,5,\dots\}$ or about solutions, if any, of the
  unrestricted equation.
\item \emph{Whether the ``good reduction at $29$'' hypotheses hold
  in any intended case.} They do not hold automatically; they are
  hypotheses of the lemma discussed in
  Section~\ref{sec:ribet}, and no Lean theorem in the repository
  shows that a hypothetical solution of \eqref{eq:gap3} must
  satisfy $29\nmid ABC$. The case $29\mid ABC$ is untouched.
\item \emph{Whether a trace mismatch eliminates the level-$32$
  newform after level lowering.} It does not, in Lean, because the
  level-lowering step itself
  (\texttt{ribet\_lowers\_Frey\_to\_level\_32}) stays an uninhabited
  \texttt{def Prop}: there is no formal link between the actual
  Galois representation of the Frey curve and the newform
  \texttt{32a1} whose trace is being compared.
\end{enumerate}
Section~\ref{sec:ribet} returns to this with the specific lemma
statement, its hypotheses, and its axiom profile.

\subsection*{The strategy, honestly labelled}

The classical strategy this development is modeled on is
Frey--Serre--Ribet--Mazur~\cite{frey1986,serre1987,ribet1990,mazur1978}:
compute the conductor of a Frey curve by Tate's
algorithm~\cite{tate1975}, exclude a rational $13$-isogeny by
Mazur's theorem on $X_0(13)$, lower the residual level in the
style of Ribet, and compare the resulting newform space with
Mordell--Weil data via the Cremona labels \texttt{26a1} and
\texttt{26b1}~\cite{cremona,lmfdb}. None of those four classical
theorems is proved, used, or even stated as a hypothesis-free
implication anywhere in the Lean files discussed below. The
Frey curve of~\eqref{eq:gap3} is the displayed Weierstrass model
\begin{equation}\label{eq:frey}
  y^2=x(x-A^4)(x+B^4)
  \qquad\text{i.e.}\qquad
  [0,\,B^4-A^4,\,0,\,-A^4B^4,\,0].
\end{equation}
The short model $[0,-(A^4+B^4),0,0,0]$ is a different curve and is
not used as $E$. Sections~\ref{sec:tate}--\ref{sec:kolyvagin} are
titled ``displayed data'' for exactly this reason: each records
what four Lean files, named \emph{after} the classical authors,
actually contain --- numerals, Finset cardinalities, and one
typeclass upgrade --- while the corresponding algebraic statements
of Tate/Mazur/Ribet/Kolyvagin remain uninhabited \texttt{def Prop}
in the parent files.

\subsection*{Pins and provenance}

The formal pin is Lean \texttt{v4.12.0}, mathlib commit
\texttt{809c3fb3b5c8f5d7dace56e200b426187516535a} with
\texttt{inputRev v4.12.0} (not \texttt{@ main}), and the
\texttt{beal-conjecture} Lake manifest prefix \texttt{83739542}.
The command \texttt{lake update} was not run. The vendor tree
\texttt{Level26/BealLevel26Foundations} at \texttt{db7a556}/\texttt{fea0c393}
has difference $0$ against \texttt{b780c9c}/\texttt{1a00fbd}.
This revision's version DOI,
\href{https://doi.org/10.5281/zenodo.22885060}{10.5281/zenodo.22885060},
was reserved through Zenodo's deposit API before publication
(Section~\ref{sec:comp}); it is a new version of \texttt{22863527}
(\texttt{v30}), which was itself a new version of \texttt{22854603}
and \texttt{22851334}; the \texttt{beal-conjecture} concept DOI
\href{https://doi.org/10.5281/zenodo.22041831}{10.5281/zenodo.22041831}
stays, related as \texttt{isVersionOf}~\cite{zenodo-v31,zenodo-concept}.
The separate \texttt{beal-level-26-foundations} repository has its
own concept DOI,
\href{https://doi.org/10.5281/zenodo.22379293}{10.5281/zenodo.22379293}~\cite{zenodo-concept-foundations};
it is not the same series as this repository's.
Its end-of-proof release
\path{v30.1.7-beal-final-43735b3} is archived as
\href{https://doi.org/10.5281/zenodo.22912430}{10.5281/zenodo.22912430},
and the post-final \path{v25.0.1-forward-43735b3} retrofit as
\href{https://doi.org/10.5281/zenodo.22922473}{10.5281/zenodo.22922473}
\cite{zenodo-foundations-v3017,zenodo-foundations-v2501}.
The resulting cross-repository integration is archived as
\href{https://doi.org/10.5281/zenodo.22923337}{10.5281/zenodo.22923337}
\cite{zenodo-integrated}. Its annotated Git tag remains on the triggering
snapshot; this follow-up on \texttt{main} adds the DOI that did not exist
when that snapshot was created.

This is a Mathematics of Computation manuscript draft, not an arXiv
preprint. There is no arXiv identifier and no arXiv badge.

The four display files live in
\texttt{Level26/HonestB0Search/} on \texttt{beal-conjecture}
\texttt{main} (release tag
\texttt{v30.0.0-level-26-structural-verification}, commit
\texttt{af520d3}) and on foundations
\texttt{phase-darmon-merel-4413} at \texttt{2b60ccd}. At this
writing the HonestB0Search slices are near identical;
\texttt{beal-conjecture} is the superset, carrying the historical
\texttt{BealMatveevBeal} target in addition to the four finals. Any
commit hash for \texttt{main} cited elsewhere in this draft for
purposes other than pinning a specific historical file is only
accurate as of the revision that introduced it; the release tag
above is the stable anchor.
